\subsection{Discrete families and close maps}\label{subsec:hilbert-preliminaries-discrete}
\paraid{p-0e8bda81dbda}
To state the approximation condition in \Cref{thm:torunczyk},
we recall discreteness of a family of subsets and closeness of maps
with respect to an open cover.

\begin{definition}\label{def:discrete-family}
\paraid{p-360939bb32bd}
A family
$\{\ClosedDomain _\FamilyIndex \}_{\FamilyIndex \in \DomainIndexSet }$
of subsets of a space
$\RecognitionTarget $
is \emph{discrete in}
$\RecognitionTarget $
if every point of
$\RecognitionTarget $
has a neighborhood meeting at most one member of the family
$\{\ClosedDomain_\FamilyIndex\}_{\FamilyIndex\in\DomainIndexSet}$.
We call the family
$\{\ClosedDomain_\FamilyIndex\}_{\FamilyIndex\in\DomainIndexSet}$
\emph{locally finite} if every point of
$\RecognitionTarget$
has a neighborhood meeting only finitely many members.
These requirements apply also at points outside the union.
For an open cover
$\OpenCover$
of
$\RecognitionTarget $,
two maps
$\InputMap ,\ApproximatingMap \colon \CommonMapDomain \to \RecognitionTarget $
are
\emph{$\OpenCover$-close}
if for each
$\DomainPoint \in \CommonMapDomain $
there exists
$U\in\OpenCover$
such that
\[
 \{\InputMap(\DomainPoint),\ApproximatingMap(\DomainPoint)\}\subset U.
\]
\end{definition}
